%basic cover letter template
\documentclass{letter}
\usepackage{amssymb,amsmath}
\usepackage{graphicx}

\oddsidemargin=.2in
\evensidemargin=.2in
\textwidth=6in
\topmargin=-1.5in
\textheight=10in


%\address{Mathematics Department\\University of Illinois\\
%1409 W. Green St\\Urbana, Illinois 61801}

\newcommand {\qed}{\mbox{$\Box$}}
\renewcommand {\iff}{\Longleftrightarrow}
\newcommand {\R}{\mathbb{R}}
\newcommand {\N}{\mathbb{N}}
\newcommand {\Q}{\mathbb{Q}}
\newcommand {\Z}{\mathbb{Z}}

\newcommand {\sub}{\mbox{SB}}

\begin{document}

\begin{letter}{\textbf{Editorial Office}\\ 
		\textbf{Physical Review Applied}}

\opening{Dear Editor,}

We submit the manuscript \textit{Wavelength-specific refractive-index calibration of air at 1762 nm for environmental phase compensation in quantum optical interfaces} by W. Wu, T. Sch\"{a}tz, and U. Warring for consideration as a Research Article in Physical Review Applied.

Free-space optical paths in trapped-ion and neutral-atom quantum platforms accumulate phase errors from environmentally driven refractive-index fluctuations. Compensating these errors requires accurate knowledge of the refractive-index sensitivity at the operating wavelength, but standard broadband refractivity formulations (Edlén, Ciddor, Mathar) lack wavelength-specific validation at atomic clock transition wavelengths relevant to trapped-ion and quantum optical platforms. The $^{138}\mathrm{Ba}^+$ clock transition at 1762 nm lies adjacent to water-vapour absorption features, suggesting that local dispersion effects may influence the humidity coefficient; no direct measurement has been reported.

Here we report an eight-month measurement campaign in which a 1762 nm probe laser referenced to a single trapped $^{138}\mathrm{Ba}^+$ ion, combined with a dual-wavelength Michelson interferometer and co-located environmental sensors, yields a wavelength-specific calibration of the refractive index of air. The measured temperature coefficient differs from the Mathar model by 0.53 \%, and the pressure coefficient by 2.08 \%, while the humidity coefficient exhibits a 17.2 \% enhancement consistent with a Kramers–Kronig analysis of nearby water-vapour resonances. For a representative 1 m free-space path under diurnal environmental swings, feed-forward compensation using these wavelength-specific coefficients reduces the residual phase error by 85 \% compared to the Mathar model, yielding a projected single-qubit gate infidelity contribution below $3\times10^{-4}$.

This manuscript is suited to Physical Review Applied because it addresses a practical bottleneck in deployed quantum optical systems with a measurement-based solution and concrete engineering deliverables. The per-meter phase-shift sensitivities reported (rad m$^{-1}$ K$^{-1}$, rad m$^{-1}$ \%RH$^{-1}$, rad m$^{-1}$ hPa$^{-1}$) can be incorporated directly into environmental phase-noise budgets of trapped-ion and neutral-atom platforms operating at this wavelength. More broadly, the methodology — pairing an atomic frequency reference with environmental monitoring to extract wavelength-specific refractivity coefficients — is transferable to other quantum-relevant wavelengths where broadband refractivity models lack validation, with implications for precision optical metrology and atmospheric sensing.

This manuscript has not been published previously and is not under consideration elsewhere. All authors have approved the submission.

Data and analysis scripts supporting the findings are openly available through Zenodo (DOI: 10.5281/zenodo.18800166).

Yours faithfully,\\

{\bf Wei Wu}\\
Doctoral Researcher, AG Sch\"{a}tz\\
Albert-Ludwigs-Universität Freiburg\\
Physikalisches Institut\\
Hermann-Herder-Str. 3\\
79104 Freiburg, Germany\\
wei.wu@physik.uni-freiburg.de
\end{letter}

\end{document}

